% 53e Olympiade Internationale de Mathématiques — Mar del Plata, Argentine.
% Énoncés officiels en français (imo-official.org). Chaque problème vaut 7 points.

\begin{center}
  \emph{Premier jour --- mardi 10 juillet 2012 --- Durée : 4 heures 30}
\end{center}

\subsection*{Problème 1}
Soit $ABC$ un triangle et $J$ le centre de son cercle exinscrit opposé au sommet $A$. Ce cercle est tangent au
côté $[BC]$ en $M$ et aux droites $(AB)$ et $(AC)$, respectivement, en $K$ et $L$. Les droites $(LM)$ et
$(BJ)$ se coupent en $F$ et les droites $(KM)$ et $(CJ)$ se coupent en $G$. Soit $S$ le point d'intersection
des droites $(AF)$ et $(BC)$ et soit $T$ le point d'intersection des droites $(AG)$ et $(BC)$.

Montrer que $M$ est le milieu du segment $[ST]$.

(Le cercle \emph{exinscrit} du triangle $ABC$ opposé au sommet $A$ est le cercle tangent au segment $[BC]$, à
la demi-droite $[AB]$ au-delà de $B$ et à la demi-droite $[AC]$ au-delà de $C$)

\subsection*{Problème 2}
Soit $n \geq 3$ un entier et soit $a_2, a_3, \ldots, a_n$ des nombres réels strictement positifs tels que
$a_2a_3\cdots a_n = 1$. Montrer que :
\[
(1 + a_2)^2(1 + a_3)^3\cdots(1 + a_n)^n \ > \ n^n.
\]

\subsection*{Problème 3}
La \emph{devinette du menteur} est un jeu joué par deux joueurs $A$ et $B$. Les règles du jeu dépendent de
deux entiers strictement positifs $k$ et $n$, connus par chacun des deux joueurs.

Au début du jeu, le joueur $A$ choisit deux entiers $x$ et $N$ vérifiant $1 \leq x \leq N$. Le joueur $A$
garde $x$ secret et, honnêtement, communique $N$ au joueur $B$. Le joueur $B$ essaye d'obtenir des
informations concernant $x$ en posant au joueur $A$ des questions comme suit : pour chaque question, $B$
choisit un ensemble arbitraire d'entiers strictement positifs $S$ (éventuellement déjà choisi pour une
question antérieure) et demande à $A$ si $x$ appartient à $S$ ; le joueur $B$ peut poser autant de telles
questions qu'il le souhaite. Après chaque question, le joueur $A$ doit immédiatement répondre par \emph{oui}
ou \emph{non}, mais il a le droit de mentir autant de fois qu'il le souhaite ; la seule restriction étant que
parmi toutes $k + 1$ réponses consécutives, au moins l'une de ces réponses doit être la vérité.

Après que $B$ ait posé autant de questions qu'il le souhaite, il doit proposer un ensemble $X$ contenant au
plus $n$ entiers strictement positifs. Si $x$ appartient à $X$, alors $B$ gagne, sinon, il perd. Montrer que :
\begin{enumerate}
  \item Si $n \geq 2^k$, alors $B$ dispose d'une stratégie gagnante.
  \item Pour tout entier suffisamment grand $k$, il existe un entier $n \geq 1.99^k$ tel que $B$ ne dispose
  pas de stratégie gagnante.
\end{enumerate}

\bigskip
\begin{center}
  \emph{Second jour --- mercredi 11 juillet 2012 --- Durée : 4 heures 30}
\end{center}

\subsection*{Problème 4}
Trouver toutes les fonctions $f : \mathbb{Z} \to \mathbb{Z}$ telles que pour tous entiers $a$, $b$, $c$
vérifiant $a + b + c = 0$, on ait l'égalité suivante :
\[
f(a)^2 + f(b)^2 + f(c)^2 \ = \ 2f(a)f(b) + 2f(b)f(c) + 2f(c)f(a).
\]
($\mathbb{Z}$ est l'ensemble des entiers relatifs)

\subsection*{Problème 5}
Soit $ABC$ un triangle dans lequel $\widehat{BCA} = 90^\circ$, et soit $D$ le pied de la hauteur issue de
$C$. Soit $X$ un point intérieur au segment $[CD]$. Soit $K$ le point du segment $[AX]$ tel que $BK = BC$. De
même, soit $L$ le point du segment $[BX]$ tel que $AL = AC$. Finalement, soit $M$ le point d'intersection des
droites $(AL)$ et $(BK)$.

Montrer que $MK = ML$.

\subsection*{Problème 6}
Trouver tous les entiers strictement positifs $n$ pour lesquels il existe des entiers positifs ou nuls
$a_1, a_2, \ldots, a_n$ tels que :
\[
\frac{1}{2^{a_1}} + \frac{1}{2^{a_2}} + \cdots + \frac{1}{2^{a_n}} \ = \ \frac{1}{3^{a_1}} + \frac{2}{3^{a_2}} + \cdots + \frac{n}{3^{a_n}} \ = \ 1.
\]
